\section{The long-return remainder}\label{sec:returns}

The sparse-pin estimate has made diffusion available under every allowed
pin law. We now improve its one-slot consequence to the order
\(\beta^{-3/2}\) needed for Bloch's law. The main step is a finite
expansion of the return probability: the first \(2K-1\) terms remain
explicit, and the expected remainder is at most
\(C K^{-1/2}\beta^{-3/2}+o_K(\beta^{-3/2})\). The constant \(C\)
must be independent of \(K\), since \(K\) will tend to infinity after
\(\beta\).

Fix an allowed pin set \(E\), a deterministic slot-time point
\(i\notin E\), and its time cut \(\theta\). Write \(x\) for its site,
\(z=\bar x\), and \(P=E\cup\{i\}\). At each exposure under
\(\mathbb P_P\), use the boundary-killed one-period kernel \(U\) and
\(Q=(U+U^*)/2\), \(A_Q=I-Q\). Boundary coordinates are retained with
value zero. Let \(\eta=\ind{\{i\}}\), and let \(f\) be the hitting
function of Lemma~\ref{lem:hitting-energy}, with target \(\{i\}\).
Thus \(f(i)=1\), \(0\le f\le1\), and \(f=U^*f\) away from \(i\).
Put \(g=f-\eta\). For the remainder of the section, \(K\ge1\) is
fixed, \(R=R_K\) and \(\tau=K\beta\) are the scales in
\eqref{eq:scales}, and every \(o_K\) is uniform in the cube, the cut,
the slot, and the number, positions, and times of the deterministic
extra pins.

\begin{lemma}[Finite return expansion]\label{lem:return-expansion}
For almost every exposure, define
\[
 v=\sum_{n=1}^{K-1}Q^n\eta,\qquad y=g-v,\qquad
 E_y=\langle y,A_Qy\rangle,
\]
with the empty sum interpreted as zero. Then \(E_y\ge0\),
\(\|y\|_\infty\le K+1\), and the fresh return probability satisfies
\begin{equation}\label{eq:return-expansion}
 q_i\le\sum_{n=1}^{2K-1}\langle\eta,Q^n\eta\rangle
          +2\langle Q^K\eta,y\rangle-E_y .
\end{equation}
\end{lemma}

\begin{proof}
Lemma~\ref{lem:hitting-energy} gives
\[
 q_i\le Q_{ii}+2\langle Q\eta,g\rangle
                  -\langle g,A_Qg\rangle.
\]
Every \(Q^n\eta\) is a nonnegative mass of total at most one.
Consequently \(\|v\|_\infty\le K-1\) and
\(\|y\|_\infty\le K\); the slightly weaker displayed bound also
covers all later uses. The contraction and symmetry of \(Q\) give
\(E_y\ge0\). The finite telescoping identity
\(A_Qv=Q\eta-Q^K\eta\) shows that substitution of \(g=v+y\)
leaves the linear term \(2\langle Q^K\eta,y\rangle\). To compute
the constant term, put \(a_n=\langle\eta,Q^n\eta\rangle\). By
symmetry,
\begin{align*}
 Q_{ii}+2\langle Q\eta,v\rangle-\langle v,A_Qv\rangle
 &=a_1+\langle v,Q\eta\rangle+\langle v,Q^K\eta\rangle\\
 &=a_1+\sum_{n=1}^{K-1}a_{n+1}+\sum_{n=1}^{K-1}a_{n+K}
   =\sum_{n=1}^{2K-1}a_n.
\end{align*}
This proves \eqref{eq:return-expansion}. No infinite Green series or
assumption of eventual boundary absorption is used.
\end{proof}

\begin{proposition}[Long-return bound]\label{prop:long-tail}
For every fixed integer \(K\ge1\), as \(\beta\to\infty\),
\begin{equation}\label{eq:long-tail}
 \mathbb E_P\bigl[2\langle Q^K\eta,y\rangle-E_y\bigr]
 \le C K^{-1/2}\beta^{-3/2}+o_K(\beta^{-3/2}).
\end{equation}
Here \(C\) depends only on the fixed interaction, spin, and geometric
choices, and is independent of \(K\). The error has the uniformity
specified above.
\end{proposition}

\begin{proof}
Set \(w=Q^{K-1}\eta\) and \(I=[\beta/3,2\beta/3]\). Then
\[
 2\langle Q^K\eta,y\rangle
 =\langle U_+w,y\rangle+\langle U_-w,y\rangle.
\]
For \(d\in\{+,-\}\) and almost every \(u\in I\), split the last
traversal as in Lemma~\ref{lem:suffix-energy}:
\[
 U_d=B_d(u)C_d(u),\qquad Y_d(u)=B_d(u)^*y.
\]
Both vectors in the resulting pairing
\[
 \langle U_dw,y\rangle=\langle C_d(u)w,Y_d(u)\rangle
\]
are on the same hole layer at physical time
\(t_d(u)=\theta+u\) or \(\theta-u\pmod\beta\).
The suffix remains killed only at the boundary. With zero extension to
the slice graph, Lemma~\ref{lem:suffix-energy} gives
\begin{equation}\label{eq:tail-suffix}
 \|Y_d(u)\|_\infty\le K+1,\qquad
 \int_I\mathcal D_{t_d(u)}(Y_d(u))\,du\le E_y.
\end{equation}
In particular the energy on the right is the original one-period
quadratic form, with no localization error.

Localize the mass in the first argument throughout its propagation
from \(i\). Namely, put
\[
 Q_R=\tfrac12(U_{+,R}+U_{-,R}),\qquad
 \mu_d(u)=C_{d,R}(u)Q_R^{K-1}\eta,
\]
where the subscript \(R\) means killing on leaving
\(\mathfrak B_R(z)\), in addition to boundary killing.
Expanding the powers of \(Q\) and \(Q_R\) gives convex combinations
of matching direction words. Lemma~\ref{lem:localization}, applied
to the whole history of each word, therefore yields
\begin{equation}\label{eq:tail-localization}
 \mathbb E_P\|C_d(u)w-\mu_d(u)\|_1\le C_K e^{-c\beta},
 \qquad u\in I.
\end{equation}

Let \(G\) denote the event that the exposure is good for \(z,R\),
and retain only regular split slices in the time average
\[
 F_d=\frac1{|I|}\int_I
  \ind{G}\ind{\{t_d(u)\text{ is regular for }z,R\}}
       \langle\mu_d(u),Y_d(u)\rangle\,du.
\]
The original pairing is independent of \(u\). Since \(\mu_d(u)\)
has mass at most one, \eqref{eq:tail-suffix},
\eqref{eq:tail-localization}, \eqref{eq:irregularity}, and
\eqref{eq:good-exposure} imply
\begin{align}
 \mathbb E_P\bigl|\langle U_dw,y\rangle-F_d\bigr|
 &\le(K+1)\left(C_K e^{-c\beta}+\mathbb P_P(G^c)
    +\frac1{|I|}\int_I
       \mathbb P_P(t_d(u)\text{ is not regular for }z,R)\,du\right)
       \notag\\
 &=o_K(\beta^{-3/2}).\label{eq:tail-discarded}
\end{align}
These are unconditional expectations under the exposure law. In
particular no independence between the propagated mass, the suffix,
and the regularity events is asserted or needed.

On a retained slice, every component of \(\mu_d(u)\) has elapsed
propagation length
\[
 b=(K-1)\beta+u\ge(K-\tfrac23)\beta\ge\tfrac13K\beta.
\]
Proposition~\ref{prop:diffusion} and convexity of the norm give
\(\|\mu_d(u)\|_2\le C\tau^{-3/4}\), with a constant independent
of \(K\). The mass is at most one and is supported over
\(\mathfrak B_R(z)\). Since \(1\le\tau\le R\), all hypotheses of
Proposition~\ref{prop:spread-energy} hold. Denote its constant for
this fixed norm bound by \(C_{\rm sp}\); it too is independent of
\(K\). Hence
\[
 |\langle\mu_d(u),Y_d(u)\rangle|
 \le\bigl(C_{\rm sp}\tau^{-1/2}
             \mathcal D_{t_d(u)}(Y_d(u))\bigr)^{1/2}.
\]
Cauchy--Schwarz in \(u\), followed by \eqref{eq:tail-suffix}, gives
\[
 |F_d|\le
 \left(\frac{3C_{\rm sp}\tau^{-1/2}}{\beta}\right)^{1/2}E_y^{1/2}.
\]
This also holds on \(G^c\), where \(F_d=0\). Completing the square
therefore gives, pointwise in the exposure,
\[
 F_++F_--E_y
 \le\frac{3C_{\rm sp}\tau^{-1/2}}{\beta}
 =3C_{\rm sp}K^{-1/2}\beta^{-3/2}.
\]
Take expectations and add the two errors in
\eqref{eq:tail-discarded}. This proves \eqref{eq:long-tail} and
identifies a leading constant that is independent of \(K\).
\end{proof}

\begin{corollary}[Uniform singleton sparsity]\label{cor:singleton-bound}
For every sufficiently large \(\beta\), uniformly in all allowed
pin sets \(E\), cubes, cuts, and slot-time points \(i\notin E\),
\begin{equation}\label{eq:singleton-return}
 r(E,\{i\})\le C\beta^{-3/2}.
\end{equation}
Consequently, for every allowed pin set \(D\) and every slot-time
point \(j\) at a deterministic time,
\begin{equation}\label{eq:free-slot}
 \mathbb P_D(j\in\mathcal B_D)\le C\beta^{-3/2}.
\end{equation}
\end{corollary}

\begin{proof}
Take \(K=1\) in Lemma~\ref{lem:return-expansion} and
Proposition~\ref{prop:long-tail}. By the fresh-walk identity,
\[
 r(E,\{i\})=\mathbb E_Pq_i
 \le\mathbb E_P Q_{ii}+C\beta^{-3/2}+o(\beta^{-3/2}).
\]
The two one-turn localization losses are at most \(Ce^{-c\beta}\)
in expectation. On a good exposure, the localized diagonal is at most
\(C\beta^{-3/2}\) by Proposition~\ref{prop:diffusion}; on its
complement it is at most one. The complement has probability
\(o(\beta^{-3/2})\) by \eqref{eq:good-exposure}. This bounds
\(\mathbb E_PQ_{ii}\) and proves \eqref{eq:singleton-return}.
For \(j\notin D\), apply this estimate with \(E=D\), using the time
of \(j\) as the cut, and use Proposition~\ref{prop:pin-identities}:
\[
 \mathbb P_D(j\in\mathcal B_D)=2d_D(j)
 =\frac{2r(D,\{j\})}{1+r(D,\{j\})}
 \le2r(D,\{j\}).
\]
For \(j\in D\), the left side is zero.
\end{proof}

The estimate \eqref{eq:free-slot} is obtained only after closing the
sparse-pin induction. It will now control encounters between exposed
strands and an independent ideal path, allowing the explicit finite
sum in \eqref{eq:return-expansion} to be evaluated sharply.
